% GENERATED FILE. Edit canonical lessons, then run npm run generate:books.
% canonical-source-hash: fnv1a64:47d57a29fd32396a
% canonical-lessons: SA-C239-asa, SA-C239-prema, SA-C239-snehah, SA-C239-daya, SA-C239-karuna

\chapter{Hope, Love, Affection, Compassion, Pity}
\label{ch:sa-hope-love-affection-compassion-pity}
\hlchaptermodality{239}

\begin{chapteropening}
\textbf{By the end of this chapter:} \emph{I can say hope, love, affection, compassion and pity.}

Complete the last lesson of chapter 239: I can say hope, love, affection, compassion and pity.
\end{chapteropening}

\section[\texorpdfstring{\sk{आशा} (āśā)}{आशा (āśā)}]{\sk{आशा} (āśā) — hope}

\label{lesson:SA-C239-asa}



\begin{quote}
Before the new one: say the Sanskrit for a policeman, then the Sanskrit for a driver.
\end{quote}

\subsection*{You'll want to know: \sk{आशा}}
\textbf{\sk{आशा}} — \emph{āśā} — \textquotedblleft{}hope\textquotedblright{}.

\begin{grammarlens}[title={What you've built}]
One more word for this chapter.
\end{grammarlens}

\subsection*{Your turn}
\begin{itemize}
\raggedright
  \item \emph{Say it:} \emph{āśā}
  \item \emph{Say it:} \emph{āśā}, once more
  \item \emph{Say it:} say \emph{āśā}
  \item \emph{Recall:} say the Sanskrit for a policeman, then the Sanskrit for a driver, then say \emph{āśā} again
\end{itemize}

\begin{tcolorbox}[breakable,colback=teal!4,colframe=teal!35!black,title={Before you move on}]
What does \sk{आशा} mean? (\textquotedblleft{}Hope\textquotedblright{}.) Say it once more.
\end{tcolorbox}
\section[\texorpdfstring{\sk{प्रेम} (prema)}{प्रेम (prema)}]{\sk{प्रेम} (prema) — love}

\label{lesson:SA-C239-prema}



\begin{quote}
Before the new one: say the Sanskrit for a driver, then the Sanskrit for hope.
\end{quote}

\subsection*{You'll want to know: \sk{प्रेम}}
\textbf{\sk{प्रेम}} — \emph{prema} — \textquotedblleft{}love\textquotedblright{}.

\begin{grammarlens}[title={What you've built}]
One more word for this chapter.
\end{grammarlens}

\subsection*{Your turn}
\begin{itemize}
\raggedright
  \item \emph{Say it:} \emph{prema}
  \item \emph{Say it:} \emph{prema}, once more
  \item \emph{Say it:} say \emph{prema}
  \item \emph{Recall:} say the Sanskrit for a driver, then the Sanskrit for hope, then say \emph{prema} again
\end{itemize}

\begin{tcolorbox}[breakable,colback=teal!4,colframe=teal!35!black,title={Before you move on}]
What does \sk{प्रेम} mean? (\textquotedblleft{}Love\textquotedblright{}.) Say it once more.
\end{tcolorbox}
\section[\texorpdfstring{\sk{स्नेहः} (snehaḥ)}{स्नेहः (snehaḥ)}]{\sk{स्नेहः} (snehaḥ) — affection}

\label{lesson:SA-C239-snehah}



\begin{quote}
Before the new one: say the Sanskrit for hope, then the Sanskrit for love.
\end{quote}

\subsection*{You'll want to know: \sk{स्नेहः}}
\textbf{\sk{स्नेहः}} — \emph{snehaḥ} — \textquotedblleft{}affection\textquotedblright{}.

\begin{grammarlens}[title={What you've built}]
One more word for this chapter.
\end{grammarlens}

\subsection*{Your turn}
\begin{itemize}
\raggedright
  \item \emph{Say it:} \emph{snehaḥ}
  \item \emph{Say it:} \emph{snehaḥ}, once more
  \item \emph{Say it:} say \emph{snehaḥ}
  \item \emph{Recall:} say the Sanskrit for hope, then the Sanskrit for love, then say \emph{snehaḥ} again
\end{itemize}

\begin{tcolorbox}[breakable,colback=teal!4,colframe=teal!35!black,title={Before you move on}]
What does \sk{स्नेहः} mean? (\textquotedblleft{}Affection\textquotedblright{}.) Say it once more.
\end{tcolorbox}
\section[\texorpdfstring{\sk{दया} (dayā)}{दया (dayā)}]{\sk{दया} (dayā) — compassion}

\label{lesson:SA-C239-daya}



\begin{quote}
Before the new one: say the Sanskrit for love, then the Sanskrit for affection.
\end{quote}

\subsection*{You'll want to know: \sk{दया}}
\textbf{\sk{दया}} — \emph{dayā} — \textquotedblleft{}compassion\textquotedblright{}.

\begin{grammarlens}[title={What you've built}]
One more word for this chapter.
\end{grammarlens}

\subsection*{Your turn}
\begin{itemize}
\raggedright
  \item \emph{Say it:} \emph{dayā}
  \item \emph{Say it:} \emph{dayā}, once more
  \item \emph{Say it:} say \emph{dayā}
  \item \emph{Recall:} say the Sanskrit for love, then the Sanskrit for affection, then say \emph{dayā} again
\end{itemize}

\begin{tcolorbox}[breakable,colback=teal!4,colframe=teal!35!black,title={Before you move on}]
What does \sk{दया} mean? (\textquotedblleft{}Compassion\textquotedblright{}.) Say it once more.
\end{tcolorbox}
\section[\texorpdfstring{\sk{करुणा} (karuṇā)}{करुणा (karuṇā)}]{\sk{करुणा} (karuṇā) — pity}

\label{lesson:SA-C239-karuna}



\begin{quote}
Before the new one: say the Sanskrit for affection, then the Sanskrit for compassion.
\end{quote}

\subsection*{You'll want to know: \sk{करुणा}}
\textbf{\sk{करुणा}} — \emph{karuṇā} — \textquotedblleft{}pity\textquotedblright{}.

\begin{grammarlens}[title={What you've built}]
One more word for this chapter.
\end{grammarlens}

\subsection*{Your turn}
\begin{itemize}
\raggedright
  \item \emph{Say it:} \emph{karuṇā}
  \item \emph{Say it:} \emph{karuṇā}, once more
  \item \emph{Say it:} say \emph{karuṇā}
  \item \emph{Recall:} say the Sanskrit for affection, then the Sanskrit for compassion, then say \emph{karuṇā} again
\end{itemize}

\begin{tcolorbox}[breakable,colback=teal!4,colframe=teal!35!black,title={Before you move on}]
What does \sk{करुणा} mean? (\textquotedblleft{}Pity\textquotedblright{}.) Say it once more.
\end{tcolorbox}
